\begin{workedexample}[Worked Example: Larmor frequency]
Nuclear spins precess at $f_0 = \gamma B_0$; for protons
$\gamma = 42.58\ \text{MHz/T}$. At $B_0 = 1.5\ \text{T}$,
\[ f_0 = 42.58\times1.5 = 63.9\ \text{MHz}; \]
at $3\ \text{T}$ it is $127.7\ \text{MHz}$. The RF coil must be tuned to this
frequency to excite and receive the spins.
\end{workedexample}

\begin{keytakeaways}
\begin{itemize}[leftmargin=1.4em,itemsep=2pt]
\item Larmor: $f_0 = \gamma B_0$; for protons $\gamma = 42.58\ \text{MHz/T}$.
\item Net magnetization aligns with $B_0$; an on-resonance $B_1$ pulse tips it into the transverse plane.
\item Higher $B_0$ raises frequency (and signal), demanding higher-frequency coils.
\end{itemize}
\end{keytakeaways}

\begin{exercises}
\item Proton Larmor frequency at $3\ \text{T}$?
\item At what field is the proton frequency $300\ \text{MHz}$?
\item Does $7\ \text{T}$ need a higher- or lower-frequency coil than $1.5\ \text{T}$?
\end{exercises}

\furtherreading{Nishimura~\cite{nishimura} (ch.~2); Haacke~\cite{haacke}.}
